% GENERATED FILE. Edit canonical lessons, then run npm run generate:books.
% canonical-source-hash: fnv1a64:a36fdc64daff4879
% canonical-lessons: ZH-C457-laojia, ZH-C457-baituo, ZH-C457-duoxie, ZH-C457-ganji-grateful, ZH-C457-fangfa, ZH-R457-review-453-457, ZH-R457-review-final-457

\chapter{Thanks for the Help}
\label{ch:zh-thanks-for-the-help}
\hlchaptermodality{457}

\begin{chapteropening}
\textbf{By the end of this chapter:} \emph{I can ask a stranger's help politely, ask a favour, say thanks a lot, say I am grateful and ask for a good way to do something.}

Complete the last lesson of chapter 457: I can ask a stranger's help politely, ask a favour, say thanks a lot, say I am grateful and ask for a good way to do something.
\end{chapteropening}

\section[láojià]{láojià — may I trouble you (asking a favour)}

\label{lesson:ZH-C457-laojia}



\begin{quote}
Before the new one: say the Mandarin for the owner of a lost item, then the Mandarin for to claim lost property.
\end{quote}

\subsection*{You'll want to know: láojià}
\emph{láojià} — \textquotedblleft{}may I trouble you (asking a favour)\textquotedblright{}.

\begin{quote}
\emph{láojià qù huǒchēzhàn zěnme zǒu} — Excuse me, how do I get to the railway station?
\end{quote}

Its characters come later; for now it is a word for the ear.

\begin{cousinweb}
\emph{láojià} is built from \emph{láo} 'to tire' and \emph{jià} 'a carriage', both inherited from Middle Chinese, \emph{law} and \emph{kaeH} (Guangyun, in Baxter's transcription), which glosses the first 'weary' and the second 'to ride'. The pair is usually explained as an old polite way of asking someone to trouble their carriage.
\end{cousinweb}

\begin{grammarlens}[title={What you've built}]
One more word for getting around.
\end{grammarlens}

\subsection*{Your turn}
\begin{itemize}
\raggedright
  \item \emph{Say it:} \emph{láojià}
  \item \emph{Say it:} \emph{láojià}, once more
  \item \emph{Say it:} \emph{láojià qù huǒchēzhàn zěnme zǒu} — Excuse me, how do I get to the railway station?
  \item \emph{Recall:} say the Mandarin for the owner of a lost item, then the Mandarin for to claim lost property, then say \emph{láojià} again
\end{itemize}

\begin{tcolorbox}[breakable,colback=teal!4,colframe=teal!35!black,title={Before you move on}]
What does láojià mean? (\textquotedblleft{}May I trouble you (asking a favour)\textquotedblright{}.) Say it once more.
\end{tcolorbox}
\section[bàituō]{bàituō — to ask a favour of someone}

\label{lesson:ZH-C457-baituo}



\begin{quote}
Before the new one: say the Mandarin for to claim lost property, then the Mandarin for may I trouble you.
\end{quote}

\subsection*{You'll want to know: bàituō}
\emph{bàituō} — \textquotedblleft{}to ask a favour of someone\textquotedblright{}.

\begin{quote}
\emph{bàituō nǐ bāng wǒ ná xíngli} — Could you carry my luggage for me, please?
\end{quote}

Its characters come later; for now it is a word for the ear.

\begin{cousinweb}
\emph{bàituō} is built from \emph{bài} 'to bow' and \emph{tuō} 'to entrust', both inherited from Middle Chinese, \emph{peajH} and \emph{thak} (Guangyun, in Baxter's transcription), which glosses the second 'to entrust'.
\end{cousinweb}

\begin{grammarlens}[title={What you've built}]
One more word for getting around.
\end{grammarlens}

\subsection*{Your turn}
\begin{itemize}
\raggedright
  \item \emph{Say it:} \emph{bàituō}
  \item \emph{Say it:} \emph{bàituō}, once more
  \item \emph{Say it:} \emph{bàituō nǐ bāng wǒ ná xíngli} — Could you carry my luggage for me, please?
  \item \emph{Recall:} say the Mandarin for to claim lost property, then the Mandarin for may I trouble you, then say \emph{bàituō} again
\end{itemize}

\begin{tcolorbox}[breakable,colback=teal!4,colframe=teal!35!black,title={Before you move on}]
What does bàituō mean? (\textquotedblleft{}To ask a favour of someone\textquotedblright{}.) Say it once more.
\end{tcolorbox}
\section[\texorpdfstring{\zh{多谢} (duōxiè)}{多谢 (duōxiè)}]{\zh{多谢} (duōxiè) — many thanks; thanks a lot}

\label{lesson:ZH-C457-duoxie}



\begin{quote}
Before the new one: say the Mandarin for may I trouble you, then the Mandarin for to ask a favour of someone.
\end{quote}

\subsection*{You'll want to know: \zh{多谢}}
\textbf{\zh{多谢}} — \emph{duōxiè} — \textquotedblleft{}many thanks; thanks a lot\textquotedblright{}.

\begin{quote}
\textbf{\zh{多谢你我找到地图了}} \emph{duōxiè nǐ wǒ zhǎodào dìtú le} — Thanks so much: I've found my map.
\end{quote}

\begin{cousinweb}
\textbf{\zh{多谢}} (\emph{duōxiè}) is built from \zh{多} 'many' and \zh{谢} 'to thank', both inherited from Middle Chinese, \emph{ta} and \emph{zjaeH} (Guangyun, in Baxter's transcription), which glosses the second 'to decline'; the sense 'to thank' is not given there.
\end{cousinweb}

\begin{grammarlens}[title={What you've built}]
One more word for getting around.
\end{grammarlens}

\subsection*{Your turn}
\begin{itemize}
\raggedright
  \item \emph{Say it:} \emph{duōxiè}
  \item \emph{Say it:} \emph{duōxiè}, once more
  \item \emph{Say it:} \emph{duōxiè nǐ wǒ zhǎodào dìtú le} — Thanks so much: I've found my map.
  \item \emph{Recall:} say the Mandarin for may I trouble you, then the Mandarin for to ask a favour of someone, then say \emph{duōxiè} again
\end{itemize}

\begin{tcolorbox}[breakable,colback=teal!4,colframe=teal!35!black,title={Before you move on}]
What does \zh{多谢} mean? (\textquotedblleft{}Many thanks; thanks a lot\textquotedblright{}.) Say it once more.
\end{tcolorbox}
\section[gǎnjī]{gǎnjī — to feel deeply grateful}

\label{lesson:ZH-C457-ganji-grateful}



\begin{quote}
Before the new one: say the Mandarin for to ask a favour of someone, then the Mandarin for thanks a lot.
\end{quote}

\subsection*{You'll want to know: gǎnjī}
\emph{gǎnjī} — \textquotedblleft{}to feel deeply grateful\textquotedblright{}.

\begin{quote}
\emph{nǐ bāng le wǒ hěn duō wǒ hěn gǎnjī} — You've helped me so much; I'm very grateful.
\end{quote}

Its characters come later; for now it is a word for the ear.

\begin{cousinweb}
\emph{gǎnjī} is built from \emph{gǎn} 'to be moved' and \emph{jī} 'to surge', both inherited from Middle Chinese, \emph{komX} and \emph{kek} (Guangyun, in Baxter's transcription), which glosses the first 'to move' and the second as a rushing wave.
\end{cousinweb}

\begin{grammarlens}[title={What you've built}]
One more word for getting around.
\end{grammarlens}

\subsection*{Your turn}
\begin{itemize}
\raggedright
  \item \emph{Say it:} \emph{gǎnjī}
  \item \emph{Say it:} \emph{gǎnjī}, once more
  \item \emph{Say it:} \emph{nǐ bāng le wǒ hěn duō wǒ hěn gǎnjī} — You've helped me so much; I'm very grateful.
  \item \emph{Recall:} say the Mandarin for to ask a favour of someone, then the Mandarin for thanks a lot, then say \emph{gǎnjī} again
\end{itemize}

\begin{tcolorbox}[breakable,colback=teal!4,colframe=teal!35!black,title={Before you move on}]
What does gǎnjī mean? (\textquotedblleft{}To feel deeply grateful\textquotedblright{}.) Say it once more.
\end{tcolorbox}
\section[\texorpdfstring{\zh{方法} (fāngfǎ)}{方法 (fāngfǎ)}]{\zh{方法} (fāngfǎ) — a way to solve a problem; means}

\label{lesson:ZH-C457-fangfa}



\begin{quote}
Before the new one: say the Mandarin for thanks a lot, then the Mandarin for to feel deeply grateful.
\end{quote}

\subsection*{You'll want to know: \zh{方法}}
\textbf{\zh{方法}} — \emph{fāngfǎ} — \textquotedblleft{}a way to solve a problem; means\textquotedblright{}.

\begin{quote}
\textbf{\zh{找不到路的时候你有什么好方法}} \emph{zhǎo bu dào lù de shíhou nǐ yǒu shénme hǎo fāngfǎ} — When you can't find the way, what do you do?
\end{quote}

\begin{cousinweb}
\textbf{\zh{方法}} (\emph{fāngfǎ}) is built from \zh{方} 'a technique' and \zh{法} 'a rule', both inherited from Middle Chinese, \emph{pjang} and \emph{pjop} (Guangyun, in Baxter's transcription), which glosses the first, among other senses, as an art or technique and the second 'a rule'.
\end{cousinweb}

\begin{grammarlens}[title={What you've built}]
One more word for getting around.
\end{grammarlens}

\subsection*{Your turn}
\begin{itemize}
\raggedright
  \item \emph{Say it:} \emph{fāngfǎ}
  \item \emph{Say it:} \emph{fāngfǎ}, once more
  \item \emph{Say it:} \emph{zhǎo bu dào lù de shíhou nǐ yǒu shénme hǎo fāngfǎ} — When you can't find the way, what do you do?
  \item \emph{Recall:} say the Mandarin for thanks a lot, then the Mandarin for to feel deeply grateful, then say \emph{fāngfǎ} again
\end{itemize}

\begin{tcolorbox}[breakable,colback=teal!4,colframe=teal!35!black,title={Before you move on}]
What does \zh{方法} mean? (\textquotedblleft{}A way to solve a problem; means\textquotedblright{}.) Say it once more.
\end{tcolorbox}
\section[(dialogue)]{Looking back: from Bags and Cases to Thanks for the Help}

\label{lesson:ZH-R457-review-453-457}



\begin{quote}
Say the Mandarin for a backpack and a travel case.
\end{quote}

\subsection*{Your turn}
\begin{itemize}
\raggedright
  \item \emph{Say it:} a backpack, a travel case, to stuff into a bag, to carry in the hand, a luggage tag
  \item \emph{Say it:} an information desk, to make an inquiry, books on the shelves, to show someone the way, a travel brochure
  \item \emph{Say it:} one's full name, one's citizenship, to be out of date, still valid, to apply for
  \item \emph{Say it:} to get separated, to turn to someone for help, a passer-by, the owner of a lost item, to claim lost property
  \item \emph{Say it:} may I trouble you, to ask a favour of someone, thanks a lot, to feel deeply grateful, a way to solve a problem
  \item \emph{Read it:} \zh{图} — each one aloud, then a word that holds it
\end{itemize}

\begin{tcolorbox}[breakable,colback=teal!4,colframe=teal!35!black,title={Before you move on}]
A backpack? (\emph{bēibāo}.) A travel case? (\emph{xínglixiāng}.) To stuff into a bag? (\emph{sāi}.) To carry in the hand? (\emph{līn}.) A luggage tag? (\emph{biāoqiān}.) An information desk? (\emph{fúwùtái}.) To make an inquiry? (\emph{zīxún}.) Books on the shelves? (\textbf{\zh{图书}}, \emph{túshū}.) To show someone the way? (\emph{zhǐlù}.) A travel brochure? (\emph{xuānchuáncè}.) One's full name? (\emph{xìngmíng}.) One's citizenship? (\emph{guójí}.) To be out of date? (\textbf{\zh{过期}}, \emph{guòqī}.) Still valid? (\emph{yǒuxiào}.) To apply for? (\emph{shēnqǐng}.) To get separated? (\emph{zǒusàn}.) To turn to someone for help? (\emph{qiúzhù}.) A passer-by? (\textbf{\zh{路人}}, \emph{lùrén}.) The owner of a lost item? (\emph{shīzhǔ}.) To claim lost property? (\emph{rènlǐng}.) May I trouble you? (\emph{láojià}.) To ask a favour of someone? (\emph{bàituō}.) Thanks a lot? (\textbf{\zh{多谢}}, \emph{duōxiè}.) To feel deeply grateful? (\emph{gǎnjī}.) A way to solve a problem? (\textbf{\zh{方法}}, \emph{fāngfǎ}.)
\end{tcolorbox}
\section[(dialogue)]{Thanks for the Help — once more}

\label{lesson:ZH-R457-review-final-457}



\begin{quote}
Say the Mandarin for may I trouble you and to ask a favour of someone.
\end{quote}

\subsection*{Your turn}
\begin{itemize}
\raggedright
  \item \emph{Say it:} may I trouble you, to ask a favour of someone, thanks a lot, to feel deeply grateful, a way to solve a problem
\end{itemize}

\begin{tcolorbox}[breakable,colback=teal!4,colframe=teal!35!black,title={Before you move on}]
May I trouble you? (\emph{láojià}.) To ask a favour of someone? (\emph{bàituō}.) Thanks a lot? (\textbf{\zh{多谢}}, \emph{duōxiè}.) To feel deeply grateful? (\emph{gǎnjī}.) A way to solve a problem? (\textbf{\zh{方法}}, \emph{fāngfǎ}.)
\end{tcolorbox}
